\chapter*{LỜI CẢM ƠN}

Trong quá trình thực hiện đồ án tốt nghiệp, em nhận thấy việc xây dựng một sản phẩm Web cần nhiều hơn giao diện và thao tác lưu dữ liệu. Với UniTrack, mỗi chức năng đều phải được xem xét theo quyền truy cập, trạng thái dữ liệu, tính nhất quán khi có thao tác đồng thời, khả năng bảo toàn minh chứng sau duyệt và trải nghiệm sử dụng của giảng viên, sinh viên.

Em xin gửi lời cảm ơn chân thành tới thầy Đỗ Quốc Huy, giảng viên hướng dẫn, đã định hướng đề tài, góp ý phạm vi sản phẩm và nhắc em giữ trọng tâm vào bài toán giám sát tiến độ dự án sinh viên. Những góp ý đó giúp em tập trung vào luồng chính: giảng viên quản lý dự án, giao nhiệm vụ, sinh viên nộp tiến độ và giảng viên phản hồi.

Em cũng xin cảm ơn các thầy cô đã trang bị cho em nền tảng về lập trình Web, cơ sở dữ liệu, phân tích thiết kế hệ thống và kiểm thử phần mềm. Trong đồ án này, các kiến thức đó được áp dụng vào nhiều phần cụ thể như thiết kế API, quản lý cookie-based session, xử lý transaction dữ liệu, xây dựng giao diện React, tổ chức dữ liệu PostgreSQL và viết kiểm thử tự động.

Cuối cùng, em xin cảm ơn gia đình và bạn bè đã động viên trong thời gian thực hiện đồ án. Sản phẩm và báo cáo vẫn còn hạn chế, nhưng quá trình hoàn thiện UniTrack đã giúp em hiểu rõ hơn cách nhìn một hệ thống phần mềm như một sản phẩm có dữ liệu, quyền truy cập, trải nghiệm người dùng và trách nhiệm bảo trì lâu dài.
